\section{第二阶段源码与实验：把任务、调度、同步落到可验证对象}

前面的章节已经把进程生命周期、调度策略和同步状态机讲成了模型。模型还不够。读者必须能把模型落到三类材料上：真实内核源码里的对象和函数、可运行的小实验、出错时能区分原因的观测命令。本节不追求把 Linux 源码逐行翻译，而是给出一条能读通的路线：先找对象，再找状态转移，再找失败路径，最后用实验验证自己的理解。

\subsection{源码阅读地图：先读结构，再读路径}

操作系统源码很大，直接从入口函数一路跳会很快迷路。更可靠的读法是先确定“哪些结构保存状态”，再读“哪些函数改变状态”。以 Linux 为参照，第二阶段最相关的路径如下。

\begin{longtable}{p{0.26\linewidth}p{0.34\linewidth}p{0.28\linewidth}}
\toprule
主题 & 典型源码位置 & 读的时候要抓住什么 \\
\midrule
任务结构 & \code{include/linux/sched.h} & \code{task\_struct} 的状态、调度实体、父子关系、引用对象 \\
fork 路径 & \code{kernel/fork.c} & \code{kernel\_clone}、\code{copy\_process}、\code{copy\_mm}、\code{copy\_files}、\code{copy\_thread} \\
exec 路径 & \code{fs/exec.c} & 二进制格式检查、参数复制、新 \code{mm} 提交点、凭据转换 \\
exit/wait & \code{kernel/exit.c} & \code{do\_exit}、zombie、reparent、wait 扫描和最终释放 \\
核心调度 & \code{kernel/sched/core.c} & enqueue/dequeue、wakeup、schedule、context switch 调用边界 \\
公平调度 & \code{kernel/sched/fair.c} & \code{vruntime}、\code{min\_vruntime}、唤醒放置、负载均衡 \\
实时调度 & \code{kernel/sched/rt.c}、\code{kernel/sched/deadline.c} & 固定优先级、预算、deadline 和准入控制 \\
futex & \code{kernel/futex/core.c} & 用户态原子值和内核等待队列如何连接 \\
mutex/rtmutex & \code{kernel/locking/mutex.c}、\code{kernel/locking/rtmutex.c} & owner、waiter、优先级继承、慢路径 \\
pipe 等待 & \code{fs/pipe.c} & 条件谓词、等待队列、关闭唤醒、非阻塞语义 \\
RCU & \code{kernel/rcu/tree.c} & 读侧临界区、宽限期、延迟释放 \\
\bottomrule
\end{longtable}

读源码时不要先追所有宏。先写出每个对象的字段表：哪些字段是身份，哪些字段是状态，哪些字段是链接关系，哪些字段是引用计数，哪些字段只在持锁时可改。然后再读改变这些字段的路径。这样能避免把源码读成函数调用列表。

\subsection{fork 的源码读法：每一步都问可见性}

\code{fork} 或 \code{clone} 的阅读重点不是“它创建了进程”，而是每一步创建出的对象什么时候对其他 CPU 可见。一个半初始化任务若提前进入 PID 表、父子列表或 runqueue，后续调度、信号、wait 都可能观察到破碎状态。

\begin{lstlisting}
source reading checklist for fork:
  1. task object and kernel stack allocated?
  2. scheduler fields initialized but not runnable?
  3. credentials, files, fs, signal, mm copied or shared?
  4. child register frame prepared?
  5. pid allocated and published?
  6. parent-child relation linked?
  7. task finally made runnable?
  8. every failure label releases exactly the acquired references?
\end{lstlisting}

这条路线能解释源码里的错误标签。若 \code{copy\_files} 成功、\code{copy\_mm} 失败，错误路径必须释放文件表引用；若 PID 已经发布，回滚就要同时撤销 PID 映射；若任务已经 runnable，很多失败不能再用简单 free，而要走退出路径。源码中的阶段边界，就是原理章节里的生命周期边界。

\subsection{上下文切换：保存的不是“整个进程”}

上下文切换常被抽象成“保存 A，恢复 B”。这句话容易误导。内核并不会每次复制整个进程资源；它保存的是继续执行所需的 CPU 上下文，并切换当前 CPU 指向的任务、栈、地址空间和调度状态。地址空间、文件表、凭据这些对象通过引用挂在 task 上，不会在每次切换时复制。

\begin{longtable}{p{0.2\linewidth}p{0.4\linewidth}p{0.28\linewidth}}
\toprule
项目 & 切换时发生什么 & 为什么重要 \\
\midrule
内核栈 & 从 prev 的内核栈切到 next 的内核栈 & 每个线程有自己的内核执行上下文 \\
通用寄存器 & 保存调用约定要求保存的部分，并恢复 next 的部分 & 保证内核函数返回后能继续执行 \\
PC/返回地址 & 由切换汇编和栈帧间接恢复 & next 像是从上次 \code{schedule} 后继续 \\
页表根 & 进程地址空间不同才需要切换或标记切换 & 影响 TLB、ASID/PCID 和内存隔离 \\
FPU/SIMD 状态 & 通常惰性或按需保存恢复 & AVX 状态很大，不能无脑复制 \\
TLS/段寄存器 & 用户线程切换时要保持线程局部状态正确 & 影响 \code{thread\_local} 和 libc 线程数据 \\
调度统计 & 记录运行时间、等待时间和切换事件 & 后续诊断依赖这些统计 \\
\bottomrule
\end{longtable}

在 x86-64 Linux 中，可以沿着 \code{schedule()}、\code{context\_switch()}、\code{switch\_to} 和架构相关切换汇编去读。读的时候重点看三件事：prev 的栈指针保存在哪里，next 的栈指针从哪里恢复，地址空间切换在哪个条件下发生。明白这三点，才算真正理解“线程切换比进程切换轻”的具体含义。

\subsection{一个可检查的调度器数据结构}

为了把调度算法写得可测，先不要急着实现所有策略。可以先写一个最小模型：任务状态、runqueue、等待原因和调度统计必须显式存在。

\begin{lstlisting}[language=C++]
#include <cstdint>
#include <deque>
#include <limits>
#include <stdexcept>
#include <vector>

using TaskId = std::size_t;
using ProcessId = std::uint32_t;

constexpr TaskId SCHED_IDLE_TASK = std::numeric_limits<TaskId>::max();
constexpr ProcessId SCHED_INVALID_PID = 0;

enum class TaskState {
  Runnable,
  Running,
  Sleeping,
  Zombie,
};

enum class WaitReason {
  None,
  Mutex,
  Io,
  Timer,
  ChildExit,
  PageFault,
};

struct Task {
  TaskId tid = SCHED_IDLE_TASK;
  ProcessId pid = SCHED_INVALID_PID;
  int priority = 0;
  TaskState state = TaskState::Runnable;
  WaitReason wait_reason = WaitReason::None;
  std::uint64_t vruntime = 0;
  std::uint64_t runnable_since = 0;
  std::uint64_t total_runtime = 0;
  std::uint64_t total_wait = 0;
};

struct RunQueue {
  std::deque<TaskId> rr;
  std::uint64_t clock = 0;
};
\end{lstlisting}

这段结构虽然小，却已经能写出硬断言：只有 \code{Runnable} 能入队；调度选中的任务必须从 \code{Runnable} 变成 \code{Running}；进入 \code{Sleeping} 时必须记录 \code{wait\_reason}；进入 \code{Zombie} 后不能再回到 runqueue；每次 tick 要更新当前任务的 runtime，而不是误把 sleeping 时间算成运行时间。

\begin{lstlisting}[language=C++]
void enqueue(RunQueue& rq, std::vector<Task>& tasks, TaskId tid) {
  auto& task = tasks.at(tid);
  if (task.state != TaskState::Runnable) {
    throw std::logic_error("only runnable task may enter runqueue");
  }
  task.runnable_since = rq.clock;
  rq.rr.push_back(tid);
}

TaskId pick_next(RunQueue& rq, std::vector<Task>& tasks) {
  while (!rq.rr.empty()) {
    const TaskId tid = rq.rr.front();
    rq.rr.pop_front();
    auto& task = tasks.at(tid);
    if (task.state == TaskState::Runnable) {
      task.total_wait += rq.clock - task.runnable_since;
      task.state = TaskState::Running;
      return tid;
    }
  }
  return SCHED_IDLE_TASK;
}
\end{lstlisting}

很多教学代码的调度器 bug 来自“队列里只有 tid，没有状态断言”。任务睡眠后忘记出队、退出后仍被选中、唤醒时重复入队，都会在压力测试下变成随机错误。把状态检查写进数据结构入口，能让错误尽早暴露。

\subsection{futex 实验：用户态快路径为什么还需要内核}

futex 的实验目标不是背 \code{futex(2)} 参数，而是观察“无竞争不进内核，有竞争才睡眠”。可以写一个用户态 mutex：先用 CAS 从 0 改到 1；失败后把状态改成 2，再调用 futex wait；解锁时把状态置 0，若旧状态是 2，则 futex wake。

\begin{lstlisting}
observe futex:
  strace -f -e futex ./futex_mutex_demo uncontended
  strace -f -e futex ./futex_mutex_demo contended
  perf stat -e context-switches,cpu-migrations ./futex_mutex_demo contended
\end{lstlisting}

无竞争版本中，绝大多数 lock/unlock 不应产生 futex 系统调用；竞争版本中，应该能看到 wait/wake。若每次加锁都进内核，说明 fast path 失败；若竞争版本 CPU 占满但 futex 调用很少，可能在忙等；若 wait 后没人 wake，要检查 unlock 是否根据旧状态唤醒。

\subsection{lost wakeup 的可复现实验}

lost wakeup 最难的地方是它看起来像偶发卡住。要让它稳定复现，可以故意写一个错误版本：等待者先检查谓词，释放锁，然后再把自己加入等待队列。唤醒者在这个窗口里改变条件并 notify，等待者随后入队睡眠，就会错过唯一一次唤醒。

\begin{lstlisting}
wrong wait sequence:
  lock(m)
  if (!ready) {
    unlock(m)              // window begins
    enqueue(waiters, self)
    sleep()
  }

waker:
  lock(m)
  ready = true
  notify_one()
  unlock(m)
\end{lstlisting}

正确实验要用 barrier 固定时序：让等待线程停在“释放锁之后、入队之前”，再让唤醒线程完成 notify，最后放行等待线程。期望结果是错误实现永久睡眠，正确实现不会睡眠或会立即重新检查谓词。这个测试比随机跑一万次更有价值，因为它直接覆盖了危险窗口。

\subsection{condition variable 的三个负面案例}

条件变量的常见错误可以拆成三个小实验，每个实验都应该有明确期望。

\begin{longtable}{p{0.25\linewidth}p{0.38\linewidth}p{0.25\linewidth}}
\toprule
错误 & 构造方法 & 期望现象 \\
\midrule
\code{if} 代替 \code{while} & 两个消费者等待一个 item，notify all 后同时醒来 & 一个消费者拿到 item，另一个必须重新睡眠；错误实现会读空队列 \\
修改状态前 notify & producer 先 notify，再 push item & consumer 醒来看到旧状态，可能继续睡或错误返回 \\
关闭不广播 & 多个 producer/consumer 睡在不同队列，close 只唤醒一边 & 另一边永久等待 \\
\bottomrule
\end{longtable}

这三个案例对应三条规则：等待必须循环检查谓词；状态改变和 notify 的顺序必须让等待者能看到新状态；关闭是一种状态变化，必须唤醒所有可能等待这个状态的线程。

\subsection{same-pool future deadlock}

线程池和 future 是现代 C++ 程序里很常见的同步组合。一个典型死锁是：线程池只有 N 个 worker，N 个任务都在等待同一个线程池中尚未运行的子任务。所有 worker 都被等待占住，子任务无人执行。

\begin{lstlisting}
bad pattern:
  worker task A:
    future = pool.submit(B)
    return future.get()   // blocks a worker

if all workers do this:
  no worker remains to run B
\end{lstlisting}

修复方向有三种。第一，把等待改成 continuation：B 完成后触发 A 的后续逻辑，不占住 worker。第二，允许 worker 在等待时帮助执行队列中的其他任务，但要小心重入和优先级。第三，把阻塞等待放到专门的 blocking pool，计算 pool 不做无限期等待。这个例子说明同步问题不只在内核；用户态运行时也必须有调度意识。

\subsection{调度观测命令：每个命令回答一个问题}

观测命令不是越多越好。关键是每个命令要回答一个具体问题。长命令不要塞进表格；实验报告里应把命令、问题和结论分开写。

\begin{lstlisting}
cat /proc/<pid>/sched
\end{lstlisting}

它回答单个任务运行了多久、等待了多久、迁移了几次。若 runnable wait 高，说明任务已经可运行但长期拿不到 CPU。

\begin{lstlisting}
perf stat -e context-switches,cpu-migrations,page-faults ./workload
\end{lstlisting}

它回答切换、迁移和缺页是否异常。上下文切换暴涨常见于线程过多、锁竞争或频繁阻塞；迁移过多可能破坏 cache 和 NUMA 局部性；缺页多时要进一步区分 minor fault、major fault 和回收。

\begin{lstlisting}
perf sched record ./workload
perf sched latency
\end{lstlisting}

这组命令回答“任务被唤醒后等了多久才真正运行”。若唤醒延迟集中在某些线程或某些 CPU 上，要继续检查 runqueue、CPU affinity、优先级和 cgroup quota。

\begin{lstlisting}
pidstat -w -p <pid> 1
\end{lstlisting}

它区分自愿上下文切换和非自愿上下文切换。自愿切换多，通常说明线程主动阻塞在 I/O、锁或等待事件；非自愿切换多，通常说明被抢占或 CPU 竞争强。

\begin{lstlisting}
cat /sys/fs/cgroup/<group>/cpu.stat
\end{lstlisting}

它回答服务组是否被 CPU 配额限流。若 throttled time 和 throttled periods 增长，p99 延迟很可能来自配额周期，而不是业务函数突然变慢。

\begin{lstlisting}
strace -f -e futex,clone,execve,wait4 ./workload
\end{lstlisting}

它回答线程、进程和 futex 系统调用行为。若大量线程停在 futex wait，就要寻找锁 owner、等待谓词和唤醒路径；若 wait4 长期不返回，要检查子进程是否退出以及父子关系是否正确。

\begin{lstlisting}
perf top
perf record ./workload
\end{lstlisting}

它们回答 CPU 真正在执行哪里。若热点在自旋、内存拷贝、哈希查找或内核软中断，优化方向完全不同。

做报告时不要只贴命令输出。要写清楚：实验负载是什么，预期瓶颈是什么，命令观察到的现象是否支持预期，若不支持下一步查什么。例如 p99 延迟高，同时 \code{cpu.stat} 显示大量 throttled time，说明要先看 CPU quota；若 \code{perf sched latency} 显示唤醒后等待长，但 cgroup 没限流，要看 runqueue、优先级和 CPU affinity。

\subsection{从源码到实验的闭环}

第二阶段的最低闭环可以按下面顺序完成。

\begin{enumerate}
\tightlist
\item 读 \code{copy\_process}，画出 fork 的阶段图和失败回滚表。
\item 写一个 fork/exec/wait 小程序，用 \code{strace -f} 验证父子返回值、exec 成功路径和 wait 回收。
\item 读 \code{schedule} 到架构切换路径，说明一次线程切换保存哪些状态。
\item 写一个 Round Robin 或公平调度模型，断言 sleeping/zombie 任务不能入队。
\item 写一个错误 condition variable 示例，用 barrier 稳定制造 lost wakeup。
\item 写一个 futex mutex 或阅读 pthread mutex 的 futex 行为，用 \code{strace -e futex} 区分快路径和慢路径。
\item 用 \code{perf sched} 或 \code{/proc/<pid>/sched} 分析一次人为制造的 runnable wait。
\end{enumerate}

完成这些实验后，读者不再只是知道概念名称，而是能把一个现象定位到对象、状态、源码路径和观测证据。到这一步，任务、调度与同步这一章才真正具备工程教材的质量。
